
Data quality is a recognized determinant of large language model performance, yet measuring it directly requires access to pretraining corpora that are rarely released. Model cards -- structured transparency documents accompanying open-weight LLM releases -- describe curation practices in human-readable prose, but whether these descriptions correlate with independently measured benchmark outcomes has not been tested at scale. This work introduces the LLM Documentation-Benchmark Registry, a dataset of n = 4,493 models linking binary curation documentation indicators extracted from HuggingFace model cards to Open LLM Leaderboard v2 benchmark scores. Using targeted family sampling to prioritize well-documented model families (LLaMA, Mistral, Qwen, Falcon, Pythia, OLMo), the registry achieves non-zero variance in three of four documentation features. A scale-and-architecture OLS baseline explains approximately 42\% of V2 benchmark variance (R-squared = 0.4247). Adding a composite documentation score yields a statistically significant but negative coefficient (beta = -3.45, p = 1.1 x $10^-8$), with a marginal R-squared increment of 0.0042. This negative association is attributed to a size-vintage confound rather than a causal effect of documentation on performance. Additionally, perplexity filtering -- a widely practiced curation technique -- is absent from all 3,749 retrieved model cards under its canonical name, indicating a vocabulary gap in current documentation standards. The registry is intended as infrastructure for future causal analysis using propensity-matched designs.
